\section{Turing Machine --- Build --- Counter Loop}
%%% generated by the blueprint pipeline from TCSlib.Complexity.TuringMachine.Build.CounterLoop
\subsection*{Turing Machine --- Build --- Counter Loop: Introduction}

Declarations of the Lean module \texttt{TCSlib.Complexity.TuringMachine.Build.CounterLoop}.

\subsection*{Turing Machine --- Build --- Counter Loop: Results}

\begin{definition}[Relabeling work symbols in an action]
\lean{Turing.Action.mapWorkSymbols}
\label{Turing.Action.mapWorkSymbols}
\leanok
\uses{Turing.Action}
Given an alphabet bijection $e$ and a one-step action, define its relabeled action by
applying $e$ to every symbol it may write on a work tape, leaving blanks unchanged and
preserving all head movements, the input-head movement, the output symbol, and the
successor state.
\end{definition}

\begin{definition}[Relabeling symbols in a configuration]
\lean{Turing.Cfg.mapWorkSymbols}
\label{Turing.Cfg.mapWorkSymbols}
\leanok
\uses{Turing.Cfg}
Given an alphabet bijection $e$ and a configuration for a fixed input, define the
relabeled configuration by applying $e$ to every symbol stored on every work tape,
leaving blank cells unchanged. The state, head positions, input, and output are
preserved.
\end{definition}

\begin{definition}[Relabeling work symbols in a machine]
\lean{Turing.MultiTapeTM.mapWorkSymbols}
\label{Turing.MultiTapeTM.mapWorkSymbols}
\leanok
\uses{Turing.Action.mapWorkSymbols, Turing.MultiTapeTM}
Given a deterministic multi-tape Turing machine with alphabet $Symbol$ and a
bijection $e$ of that alphabet, define a machine over the same alphabet and state set
with the same initial state. Its transition first translates each scanned work-tape
symbol back along $e$, applies the original transition, and then relabels every symbol
that the action writes on a work tape by $e$; all other parts of the action are
preserved.
\end{definition}

\begin{lemma}[Work-symbol relabelling preserves a step]
\lean{Turing.counter_work_step}
\label{Turing.counter_work_step}
\leanok
\uses{Turing.Action.apply, Turing.Action.mapWorkSymbols, Turing.Cfg, Turing.Cfg.inputSymbol, Turing.Cfg.mapWorkSymbols, Turing.Cfg.workTapeSymbols, Turing.MultiTapeTM, Turing.MultiTapeTM.mapWorkSymbols, Turing.MultiTapeTM.step}
\difficulty{5}
Let $k$ be a natural number, let $M$ be a $k$-work-tape machine over a symbol type
$\mathrm{Symbol}$ with state type $\mathrm{State}$, and let
$e:\mathrm{Symbol}\simeq\mathrm{Symbol}$ be a bijection. For any finite input word over
$\mathrm{Symbol}$ and any configuration $d$, taking one step after relabelling the
work-tape symbols by $e$ gives the same configuration as taking a step first and then
relabelling the work-tape symbols by $e$.
\end{lemma}

\begin{theorem}[Work-symbol relabelling commutes with a run]
\lean{Turing.MultiTapeTM.mapWorkSymbols_runFrom}
\label{Turing.MultiTapeTM.mapWorkSymbols_runFrom}
\leanok
\uses{Turing.Cfg, Turing.Cfg.mapWorkSymbols, Turing.MultiTapeTM, Turing.MultiTapeTM.mapWorkSymbols, Turing.MultiTapeTM.runFrom, Turing.MultiTapeTM.runFrom_comm_of_step, Turing.counter_work_step}
\difficulty{1}
Let $M$ be a multi-tape Turing machine with symbol set $S$, and let $e:S\simeq S$ be a
bijection. For any input word $x$, configuration $c$ on $x$, and number of steps
$t\in\mathbb{N}$, the configuration reached in $t$ steps by the machine whose
work-tape symbols are relabelled by $e$, starting from $c$ with its work-tape symbols
relabelled by $e$, is the configuration reached by $M$ from $c$ in $t$ steps with its
work-tape symbols relabelled by $e$.
\end{theorem}

\begin{definition}[Fixed-width Boolean decrement]
\lean{Turing.decFixed}
\label{Turing.decFixed}
\leanok
Define a partial operation on finite Boolean words as follows: it is undefined on the
empty word; on a word whose first bit is true, it changes that bit to false and leaves
the rest unchanged; on a word whose first bit is false, it recursively decrements the
remaining word and, if that succeeds, changes the first bit to true.
\end{definition}

\begin{lemma}[Decrement and complemented increment]
\lean{Turing.counter_dec_complement}
\label{Turing.counter_dec_complement}
\leanok
\uses{Turing.decFixed, Turing.incFixed}
\difficulty{4}
For every Boolean word $w$, decrementing $w$ at fixed width gives the same optional word
as first complementing every bit of $w$, incrementing at that width, and then
complementing every bit of the result; failure due to overflow is preserved.
\end{lemma}

\begin{lemma}[Involutive work-tape complementation]
\lean{Turing.counter_complement_twice}
\label{Turing.counter_complement_twice}
\leanok
\uses{Turing.Cfg, Turing.Cfg.mapWorkSymbols}
\difficulty{3}
For any configuration $d$ with $k$ work tapes over the Boolean alphabet, complementing
every work-tape symbol twice returns exactly $d$.
\end{lemma}

\begin{lemma}[Complementing a buffered word]
\lean{Turing.counter_buffer_complement}
\label{Turing.counter_buffer_complement}
\leanok
\uses{Turing.FinTM.bufferTape, Turing.FinTM.bufferTape_map}
\difficulty{1}
For every Boolean word $w$ and integer cell position $p$, the tape storing the bitwise
complement of $w$ at $p$ contains the complement of the symbol stored by the tape for
$w$ at $p$, with a blank remaining blank.
\end{lemma}

\begin{definition}[Boolean decrement machine]
\lean{Turing.decrementTM}
\label{Turing.decrementTM}
\leanok
\uses{Turing.FlagPhase, Turing.MultiTapeTM, Turing.MultiTapeTM.mapWorkSymbols, Turing.incrementTM}
For a natural number $k$ and a chosen work-tape index $i<k$, define the decrement
machine to be the increment machine for that tape with Boolean work symbols relabeled by
complementation.
\end{definition}

\begin{lemma}[Decrement run under symbol complementation]
\lean{Turing.counter_decrement_run}
\label{Turing.counter_decrement_run}
\leanok
\uses{Turing.Cfg, Turing.Cfg.mapWorkSymbols, Turing.FlagPhase, Turing.MultiTapeTM.mapWorkSymbols_runFrom, Turing.MultiTapeTM.runFrom, Turing.counter_complement_twice, Turing.decrementTM, Turing.incrementTM}
\difficulty{1}
Let $k$ be a natural number, let $x$ be a word, let $i$ be one of the $k$ work-tape
indices, let $d$ be a configuration of the increment machine on input $x$, and let $t$
be a natural number. Running the decrement machine for $t$ steps from $d$ gives the same
configuration as complementing all work-tape symbols in $d$, running the increment
machine for $t$ steps, and complementing all work-tape symbols in the result.
\end{lemma}

\begin{theorem}[Decrement machine on a non-underflowing word]
\lean{Turing.decrementTM_run_succ_ofCfg}
\label{Turing.decrementTM_run_succ_ofCfg}
\leanok
\uses{Turing.Cfg, Turing.Cfg.mapWorkSymbols, Turing.FinTM.bufferTape, Turing.FlagPhase, Turing.MultiTapeTM.runFrom, Turing.counter_buffer_complement, Turing.counter_dec_complement, Turing.counter_decrement_run, Turing.decFixed, Turing.decrementTM, Turing.incFixed, Turing.incrementTM_run_succ_ofCfg}
\difficulty{5}
Let $k$ be a number of work tapes, choose one tape $i$, and let $x$ be an input word and
$w,v$ Boolean words. Suppose fixed-width decrement sends $w$ to $v$, the initial
configuration $d$ is in the run phase, and tape $i$ agrees with the buffer storing $w$
at every offset $p$ from its head position with $-1\le p\le |w|$. After $2|(w\text{
begins with false})|+2$ steps, the decrement machine is in the done phase with flag
true; on tape $i$, the cells from the initial head position through the next $|w|-1$
cells contain the buffered word $v$, and all other tape contents are unchanged. It has
not reached either done phase at any earlier time, and through that time tape $i$'s head
stays between one cell left of its initial position and the initial position plus the
number of leading false bits of $w$, inclusive; every other tape head remains fixed.
\end{theorem}

\begin{theorem}[Decrement machine on underflow]
\lean{Turing.decrementTM_run_underflow_ofCfg}
\label{Turing.decrementTM_run_underflow_ofCfg}
\leanok
\uses{Turing.Cfg, Turing.Cfg.mapWorkSymbols, Turing.FinTM.bufferTape, Turing.FlagPhase, Turing.MultiTapeTM.runFrom, Turing.counter_buffer_complement, Turing.counter_dec_complement, Turing.counter_decrement_run, Turing.decFixed, Turing.decrementTM, Turing.incFixed, Turing.incrementTM_run_overflow_ofCfg}
\difficulty{5}
Let $k$ be a number of work tapes, choose one tape $i$, and let $x$ be an input word and
$w$ a Boolean word. Suppose fixed-width decrement of $w$ underflows, the initial
configuration $d$ is in the run phase, and tape $i$ agrees with the buffer storing $w$
at every offset $p$ from its head position with $-1\le p\le |w|$. After $2|w|+2$ steps,
the decrement machine is in the done phase with flag false; on tape $i$, the cells from
the initial head position through the next $|w|-1$ cells contain $|w|$ true bits, and
all other tape contents are unchanged. It has not reached either done phase at any
earlier time, and through that time tape $i$'s head stays between one cell left of its
initial position and the initial position plus $|w|$, inclusive; every other tape head
remains fixed.
\end{theorem}

\begin{theorem}[Input independence of the decrement transition]
\lean{Turing.decrementTM_tr_input}
\label{Turing.decrementTM_tr_input}
\leanok
\uses{Turing.FlagPhase, Turing.decrementTM}
\difficulty{1}
For any number $k$ of work tapes, tape index $i$, phase $q$, and assignment of possibly
absent Boolean symbols to the work-tape heads, the decrement machine's transition action
is the same for every choice of optional Boolean input symbol.
\end{theorem}

\begin{theorem}[Transition action of the decrement machine]
\lean{Turing.decrementTM_tr_action}
\label{Turing.decrementTM_tr_action}
\leanok
\uses{Turing.Action.mapWorkSymbols, Turing.FlagPhase, Turing.MultiTapeTM.mapWorkSymbols, Turing.decrementTM, Turing.incrementTM}
\difficulty{4}
For any natural number $k$, any work-tape index $i\in\mathrm{Fin}(k)$, any phase $q$ of the machine,
any optional Boolean input symbol, and any assignment of an optional Boolean symbol to
each work tape, the decrement machine's transition action leaves the input head
stationary, and on every work tape its write instruction is never the instruction to
write a blank.
\end{theorem}

\begin{definition}[Counter word after repeated decrements]
\lean{Turing.counterWord}
\label{Turing.counterWord}
\leanok
\uses{Turing.decFixed}
For a finite Boolean word $w$ and a natural number $r$, define its counter word after
$r$ steps by iterating $r$ times the operation that decrements the word when the
decrement is defined and otherwise leaves it unchanged.
\end{definition}

\begin{theorem}[Zero counter iterations]
\lean{Turing.counterWord_zero}
\label{Turing.counterWord_zero}
\leanok
\uses{Turing.counterWord}
\difficulty{1}
For every Boolean word $w$, applying the counter-word operation zero times returns $w$.
\end{theorem}

\begin{definition}[Counter decrement overhead]
\lean{Turing.counterOverhead}
\label{Turing.counterOverhead}
\leanok
\uses{Turing.counterWord}
For a finite Boolean word $w$ and a natural number $r$, define its counter overhead to
be the sum over $j=0,\ldots,r-1$ of twice the length of the initial run of false
bits in the counter word after $j$ steps, plus two.
\end{definition}

\begin{theorem}[Zero counter overhead]
\lean{Turing.counterOverhead_zero}
\label{Turing.counterOverhead_zero}
\leanok
\uses{Turing.counterOverhead}
\difficulty{1}
For every Boolean word $w$, the counter overhead accumulated over zero rounds is zero.
\end{theorem}

\begin{theorem}[Successive counter overhead]
\lean{Turing.counterOverhead_succ}
\label{Turing.counterOverhead_succ}
\leanok
\uses{Turing.counterOverhead, Turing.counterWord}
\difficulty{1}
For every Boolean word $w$ and natural number $r$, the counter overhead through round
$r+1$ equals the overhead through round $r$ plus $2\ell_r+2$, where $\ell_r$ is the
number of consecutive false bits at the beginning of the counter word after $r$
decrement steps.
\end{theorem}

\begin{lemma}[Binary decrement specification]
\lean{Turing.counter_dec_spec}
\label{Turing.counter_dec_spec}
\leanok
\uses{Turing.decFixed}
\difficulty{5}
For every Boolean word $w$, if the decrement operation is undefined, then $w$ consists
entirely of false entries. If it returns a word $v$, then $v$ has the same length as
$w$, the binary value of $w$ is one more than the binary value of $v$, and $v$ has one
fewer true entry than $w$ plus the length of the initial run of false entries in $w$.
\end{lemma}

\begin{lemma}[Successive counter words]
\lean{Turing.counter_word_succ}
\label{Turing.counter_word_succ}
\leanok
\uses{Turing.counterWord, Turing.decFixed}
\difficulty{1}
For every finite Boolean word $w$ and natural number $r$, the counter word after $r+1$
iterations is the result of applying one fixed-width decrement to the counter word after
$r$ iterations, using that word itself if the decrement is undefined.
\end{lemma}

\begin{lemma}[Zero value and undefined decrement]
\lean{Turing.counter_dec_zero}
\label{Turing.counter_dec_zero}
\leanok
\uses{Turing.counter_dec_spec, Turing.decFixed}
\difficulty{4}
For every Boolean word $w$, the decrement operation is undefined exactly when the binary
value of $w$ is zero.
\end{lemma}

\begin{theorem}[Length preservation under counter steps]
\lean{Turing.counterWord_length}
\label{Turing.counterWord_length}
\leanok
\uses{Turing.counterWord, Turing.counter_dec_spec, Turing.counter_word_succ, Turing.decFixed}
\difficulty{4}
For every Boolean word $w$ and natural number $r$, the counter word obtained after $r$
decrement steps has the same length as $w$.
\end{theorem}

\begin{theorem}[Value after counter decrements]
\lean{Turing.counterWord_value}
\label{Turing.counterWord_value}
\leanok
\uses{Turing.counterWord, Turing.counter_dec_spec, Turing.counter_dec_zero, Turing.counter_word_succ, Turing.decFixed}
\difficulty{5}
For every Boolean word $w$ and natural number $r$ with $r\le V(w)$, where $V(w)$ is the
binary value of $w$, the counter word after $r$ decrement steps has binary value
$V(w)-r$.
\end{theorem}

\begin{lemma}[Decrementing a counter word]
\lean{Turing.counter_word_decrement}
\label{Turing.counter_word_decrement}
\leanok
\uses{Turing.counterWord, Turing.counterWord_value, Turing.counter_dec_zero, Turing.counter_word_succ, Turing.decFixed}
\difficulty{4}
Let $w$ be a finite Boolean word and let $r$ be a natural number strictly less than the
number represented by the binary digits of $w$. Then the fixed-width decrement of the
word obtained after $r$ counter iterations is defined and equals the word obtained after
$r+1$ counter iterations.
\end{lemma}

\begin{lemma}[Counter potential identity]
\lean{Turing.counter_potential}
\label{Turing.counter_potential}
\leanok
\uses{Turing.counterOverhead, Turing.counterWord, Turing.counter_dec_spec, Turing.counter_word_decrement}
\difficulty{4}
Let $w$ be a Boolean word and let $r$ be a natural number no greater than the binary
value of $w$. Then the counter overhead after $r$ iterations plus twice the number of
true entries in $w$ equals four times $r$ plus twice the number of true entries in the
counter word after $r$ iterations.
\end{lemma}

\begin{theorem}[Counter overhead bound for an initial segment]
\lean{Turing.counterOverhead_le_of_le}
\label{Turing.counterOverhead_le_of_le}
\leanok
\uses{Turing.counterOverhead, Turing.counterWord, Turing.counterWord_length, Turing.counter_potential}
\difficulty{3}
For every Boolean word $w$ and natural number $r$ with $r\le V(w)$, where $V(w)$ is the
binary value of $w$, the counter overhead accumulated over the first $r$ rounds is at
most $4r+2|w|$.
\end{theorem}

\begin{lemma}[Exhausting a counter word]
\lean{Turing.counter_word_exhausted}
\label{Turing.counter_word_exhausted}
\leanok
\uses{Turing.counterWord, Turing.counterWord_length, Turing.counterWord_value, Turing.counter_dec_spec, Turing.counter_dec_zero}
\difficulty{3}
For every finite Boolean word $w$, iterating the fixed-width counter operation as many
times as the number represented by its binary digits yields the all-false word of the
same length.
\end{lemma}

\begin{theorem}[Bound on counter overhead]
\lean{Turing.counterOverhead_le}
\label{Turing.counterOverhead_le}
\leanok
\uses{Turing.counterOverhead, Turing.counter_potential, Turing.counter_word_exhausted}
\difficulty{4}
For every Boolean word $w$, the counter overhead accumulated over the first $V(w)+1$
rounds, where $V(w)$ is the binary value of $w$, is at most $4V(w)+2|w|+2$.
\end{theorem}

\begin{definition}[Counter loop states]
\lean{Turing.CounterLoopState}
\label{Turing.CounterLoopState}
\leanok
\uses{Turing.FlagPhase}
For a type $S$, define a counter-loop state to be one of: a body-running state carrying
an element of $S$; a decrement-running state carrying a flag phase; a designated exit
state reached after counter underflow; or a designated escape state reached when the
body reaches its specified exit.
\end{definition}

\begin{definition}[Counter loop state redirection]
\lean{Turing.counterLoopRedirect}
\label{Turing.counterLoopRedirect}
\leanok
\uses{Turing.CounterLoopState}
Given a state set $S$, a distinguished state $a\in S$, an optional designated exit
state, and a state $s\in S$, redirect $s$ to the decrement phase in its run mode if
$s=a$; otherwise redirect it to the escape state if it is the designated exit, and to
the body state carrying $s$ in all other cases.
\end{definition}

\begin{definition}[Counter decrement exit routing]
\lean{Turing.counterLoopDecExit}
\label{Turing.counterLoopDecExit}
\leanok
\uses{Turing.CounterLoopState, Turing.FlagPhase}
Given a state set $S$, a distinguished state $a\in S$, and a flag phase, map the
completed phase with flag true to the body state carrying $a$, and the completed phase
with flag false to the done state. Map every other phase to the decrement state carrying
that phase.
\end{definition}

\begin{definition}[Counter loop machine]
\lean{Turing.counterLoopTM}
\label{Turing.counterLoopTM}
\leanok
\uses{Turing.Action.mapState, Turing.CounterLoopState, Turing.MultiTapeTM, Turing.counterLoopDecExit, Turing.counterLoopRedirect, Turing.decrementTM, Turing.embedEmitTM}
Given a deterministic machine $M$ with $k$ work tapes, Boolean work symbols, and state
set $S$, a distinguished state $a\in S$, and an optional designated exit state, form a
machine with $k+1$ work tapes and states for the body, counter decrement, completion,
and escape. It starts in the decrement phase in run mode; in body states it follows $M$
on the first $k$ tapes and redirects successor states according to $a$ and the
designated exit, in decrement states it runs the counter-decrement routine on the added
tape and routes its exits, and in the done and escape states it stays in that state
without moving or writing.
\end{definition}

\begin{definition}[Embedding into the final coordinate]
\lean{Turing.lastEmb}
\label{Turing.lastEmb}
\leanok
For each natural number $k$, the embedding from the one-element finite set into the
finite set of size $k+1$ sends its unique element to the last coordinate.
\end{definition}

\begin{theorem}[Last-position embedding]
\lean{Turing.lastEmb_apply}
\label{Turing.lastEmb_apply}
\leanok
\uses{Turing.lastEmb}
\difficulty{1}
For every natural number $k$ and every element $i$ of the singleton index set, the
embedding from that singleton into the set of $k+1$ indices sends $i$ to the final index
$k$.
\end{theorem}

\begin{definition}[Counter loop over an embedded alphabet]
\lean{Turing.counterLoopTMVia}
\label{Turing.counterLoopTMVia}
\leanok
\uses{Turing.Action.mapState, Turing.CounterLoopState, Turing.MultiTapeTM, Turing.MultiTapeTM.liftSymbols, Turing.counterLoopDecExit, Turing.counterLoopRedirect, Turing.decrementTM, Turing.embedEmitTM, Turing.lastEmb}
Given a symbol type $\Gamma$, a state set $S$, a Boolean-symbol machine $M$ with $k$
work tapes and state set $S$, an embedding of Boolean symbols into $\Gamma$, a
distinguished state $a\in S$, and an optional designated exit state, form a machine over
$\Gamma$ with $k+1$ work tapes and the counter-loop states. It starts in the decrement
phase in run mode; body states follow $M$ on the first $k$ tapes, while decrement states
run the one-tape decrement routine with its Boolean symbols embedded into $\Gamma$ on
the added tape, with both phases routing successor states according to the counter-loop
rules; done and escape states stay fixed without moving or writing.
\end{definition}

\begin{definition}[Counter loop configuration]
\lean{Turing.counterLoopCfg}
\label{Turing.counterLoopCfg}
\leanok
\uses{Turing.Cfg, Turing.Cfg.mapState, Turing.CounterLoopState, Turing.embedEmitCfg}
Given a machine configuration $c$ on input $x$ with $k$ work tapes, an optional
counter-loop state $q$, a counter-tape content function
$ct:\mathbb{Z}\to\mathrm{Option}(\mathrm{Symbol})$, and a counter-head position
$cp\in\mathbb{Z}$, form a configuration on the same input with $k+1$ work tapes. Its
first $k$ tapes, their head positions, the input-head position, and output are those of
$c$, its added tape and head position are $ct$ and $cp$, and its state is $q$.
\end{definition}

\begin{definition}[Counter loop round orbit]
\lean{Turing.counterLoopOrbit}
\label{Turing.counterLoopOrbit}
\leanok
\uses{Turing.Cfg, Turing.MultiTapeTM, Turing.roundOrbit}
For a machine $M$ with $k$ work tapes, a function $\tau$ assigning a natural-number time
bound to each configuration, an initial configuration $c$ on input $x$, and a natural
number $r$, the counter loop orbit at $r$ is the round orbit of $M$ from $c$ under
$\tau$ at round $r$.
\end{definition}

\begin{lemma}[Counter-loop configuration layout]
\lean{Turing.counter_layout}
\label{Turing.counter_layout}
\leanok
\uses{Turing.Cfg, Turing.Cfg.mapState, Turing.CounterLoopState, Turing.MultiTapeTM, Turing.MultiTapeTM.runFrom_zero, Turing.counterLoopCfg, Turing.embedEmitCfg_selected_pos, Turing.embedEmitCfg_selected_tape, Turing.embedEmitTM_frame}
\difficulty{4}
Let $k$ be a natural number, let $x$ be an input word, let $q$ be an optional
counter-loop state, and let $c$ be a configuration with $k$ work tapes on input $x$.
Given counter-tape contents $ct$ indexed by integers and a counter-head position $cp$,
the counter-loop configuration has state $q$, the same input-head position and output as
$c$, and work tapes and head positions obtained by appending $ct$ and $cp$ as the final
tape and head position.
\end{lemma}

\begin{theorem}[Counter configuration components]
\lean{Turing.counter_cfg_eq}
\label{Turing.counter_cfg_eq}
\leanok
\uses{Turing.Cfg, Turing.CounterLoopState, Turing.counterLoopCfg, Turing.counter_layout}
\difficulty{1}
Let $c$ be a configuration with $k$ work tapes on input $x$, let $q$ be an optional
counter-loop state, and let $ct:\mathbb{Z}\to\mathrm{Option}(\mathrm{Bool})$ and
$cp\in\mathbb{Z}$. The configuration with $k+1$ work tapes formed by adjoining $ct$ as
the last tape, placing its head at $cp$, and setting the state to $q$ has the same
input-head position and output as $c$, and its first $k$ work tapes and head positions
are those of $c$.
\end{theorem}

\begin{lemma}[Injectivity of the decrement exit map]
\lean{Turing.counter_dec_injective}
\label{Turing.counter_dec_injective}
\leanok
\uses{Turing.counterLoopDecExit}
\difficulty{3}
For every state type $S$ and designated state $a\in S$, the map from flag phases to
counter-loop states that sends the successful completed phase to the body state at $a$,
the unsuccessful completed phase to the finished state, and every other phase to the
corresponding decrement phase is injective.
\end{lemma}

\begin{lemma}[Injectivity of the counter-loop redirect]
\lean{Turing.counter_redirect_injective}
\label{Turing.counter_redirect_injective}
\leanok
\uses{Turing.counterLoopRedirect}
\difficulty{4}
Let $S$ be a type with decidable equality, choose a state $a \in S$, and choose an
optional exit state. The map sending each $s \in S$ to the decrement-running state when
$s=a$, to the escape state when the chosen exit equals $s$, and otherwise to the body
state carrying $s$, is injective.
\end{lemma}

\begin{definition}[Counter loop execution span]
\lean{Turing.counterSpan}
\label{Turing.counterSpan}
\leanok
\uses{Turing.Cfg, Turing.CounterLoopState, Turing.MultiTapeTM, Turing.MultiTapeTM.runFrom}
Let $k$ be a natural number, let $\Gamma$ be a symbol type, and let $S$ be a state type.
Let $H$ be a machine with $k+1$ work tapes and counter-loop states, let $z$ be a
configuration on input $x$, let $t$ and $width$ be natural numbers, let $cp$ be an
integer, and let $R$ be a predicate on the head positions of the first $k$ work tapes.
Require that at every time $u<t$, the state reached from $z$ after $u$ steps is neither
done nor escape; at every time $u\le t$, the added tape's head position lies in the
integer interval from $cp-1$ to $cp+width$, inclusive, and the first $k$ work-tape head
positions satisfy $R$.
\end{definition}

\begin{lemma}[Concatenating counter spans]
\lean{Turing.counter_span_append}
\label{Turing.counter_span_append}
\leanok
\uses{Turing.Cfg, Turing.CounterLoopState, Turing.MultiTapeTM, Turing.MultiTapeTM.runFrom, Turing.MultiTapeTM.runFrom_add, Turing.counterSpan}
\difficulty{5}
Let $k$ be a natural number, let $\Gamma$ and $S$ be types, let $x$ be a finite input
word over $\Gamma$, and let $H$ be a machine with $k+1$ work tapes and states of the
counter-loop form over $S$. For an initial configuration $z$, natural times $s,t$, an
integer counter position $cp$, a natural width, and predicates $R_1,R_2,R$ on the
positions of the first $k$ work heads, suppose the run from $z$ spans time $s$ with
invariant $R_1$, and the run from its time-$s$ configuration spans time $t$ with
invariant $R_2$. If each of $R_1$ and $R_2$ implies $R$, then the run from $z$ spans
time $s+t$ with invariant $R$: it avoids both exit anchors before that time, keeps the
final work-head position in the interval from $cp-1$ to $cp+\mathrm{width}$ at every
time up to it, and satisfies $R$ for the first $k$ head positions throughout.
\end{lemma}

\begin{definition}[Counter tape patch]
\lean{Turing.counterPatch}
\label{Turing.counterPatch}
\leanok
\uses{Turing.FinTM.bufferTape}
Given a tape $c\colon\mathbb{Z}\to\mathrm{Option}(\Gamma)$, a starting position
$p\in\mathbb{Z}$, and a finite word $w$ over $\Gamma$, define a new tape by placing the
buffer encoding of $w$ at positions $p$ through $p+|w|-1$, with each position indexed
relative to $p$, and leaving all other tape cells unchanged.
\end{definition}

\begin{lemma}[Idempotence of counter patching]
\lean{Turing.counter_patch_twice}
\label{Turing.counter_patch_twice}
\leanok
\uses{Turing.counterPatch}
\difficulty{3}
Let $ct$ be an integer-indexed tape of optional symbols, let $cp$ be an integer, and let
$u$ and $v$ be words over an alphabet. If $u$ and $v$ have equal lengths, then patching
$ct$ with $u$ at $cp$ and then patching the result with $v$ at the same position gives
the same tape as patching $ct$ directly with $v$ at $cp$.
\end{lemma}

\begin{lemma}[Counter patch on the frame]
\lean{Turing.counter_patch_frame}
\label{Turing.counter_patch_frame}
\leanok
\uses{Turing.FinTM.bufferTape, Turing.FinTM.bufferTape_left, Turing.FinTM.bufferTape_nat, Turing.counterPatch}
\difficulty{4}
Let $w$ and $v$ be words over an alphabet, let $ct$ be an integer-indexed tape of
optional symbols, and let $cp$ be an integer. Suppose $v$ and $w$ have equal lengths,
and throughout the range $-1 \le p \le |w|$, the tape $ct$ at position $cp+p$ agrees
with the buffered tape storing $w$ at position $p$. Then throughout the range $-1 \le p
\le |v|$, patching $ct$ with $v$ beginning at $cp$ yields the buffered tape storing $v$
at position $cp+p$.
\end{lemma}

\begin{lemma}[Counter decrement and tape preservation]
\lean{Turing.counter_debit}
\label{Turing.counter_debit}
\leanok
\uses{Turing.Action.mapWorkSymbols, Turing.Cfg, Turing.Cfg.decodeSymbols, Turing.Cfg.liftOver, Turing.Cfg.liftOver_decodeSymbols, Turing.Cfg.liftSymbols, Turing.Cfg.mapState, Turing.CounterLoopState, Turing.FinTM.bufferTape, Turing.FinTM.bufferTape_map, Turing.FinTM.bufferTape_map_decode, Turing.FlagPhase, Turing.MultiTapeTM, Turing.MultiTapeTM.liftSymbols, Turing.MultiTapeTM.liftSymbols_runFrom_over, Turing.MultiTapeTM.mapWorkSymbols, Turing.MultiTapeTM.runFrom, Turing.MultiTapeTM.runFrom_mapState_of_agreeOn, Turing.MultiTapeTM.runFrom_zero, Turing.counterLoopCfg, Turing.counterLoopDecExit, Turing.counterLoopTMVia, Turing.counterPatch, Turing.counterSpan, Turing.counter_dec_injective, Turing.counter_dec_spec, Turing.counter_layout, Turing.decFixed, Turing.decodeBool, Turing.decodeBool_eq_some, Turing.decrementTM, Turing.decrementTM_run_succ_ofCfg, Turing.decrementTM_run_underflow_ofCfg, Turing.embedEmitCfg, Turing.embedEmitCfg_selected_pos, Turing.embedEmitCfg_selected_tape, Turing.embedEmitTM, Turing.embedEmitTM_frame, Turing.embedEmitTM_runFrom, Turing.incrementTM, Turing.lastEmb}
\difficulty{7}
Let $k$ be a natural number, let $S$ and $\Gamma$ be types with decidable equality, let
$e$ be an embedding of the Boolean alphabet into $\Gamma$, and let $M$ be a
$k$-work-tape machine over $\Gamma$ with state set $S$. Choose a designated state $a\in
S$, an optional exit state, a configuration $c$ on input $x$, a Boolean word $w$, and a
counter tape $ct:\mathbb{Z}\to\operatorname{Option}(\Gamma)$ whose head is at
$cp\in\mathbb{Z}$. Assume that for every integer $p$ with $-1\le p\le |w|$, the content
of $ct$ at $cp+p$ equals the content at $p$ of the buffer storing the word obtained by
applying $e$ to each symbol of $w$. Let $r$ be the number of initial false symbols of
$w$, and let $v$ be the fixed-length decrement of $w$, with the all-true word of length
$|w|$ used if that decrement is undefined. After $2r+2$ steps of the counter-loop
machine, started from the decrement phase with configuration $c$ and counter tape $ct$,
the configuration is the one with state equal to the body state at $a$ if the decrement
is defined and the finished state otherwise; its first $k$ work tapes and head positions
are those of $c$, and its counter tape is $ct$ patched on the interval from $cp$ through
$cp+|w|-1$ with the buffer storing the image of $v$ under $e$. At every earlier time,
the machine is neither finished nor escaped; through time $2r+2$, the counter head stays
in $[cp-1,cp+|w|]$, and the first $k$ work-head positions remain those of $c$.
\end{lemma}

\begin{lemma}[Simulation of a counter-loop body]
\lean{Turing.counter_body}
\label{Turing.counter_body}
\leanok
\uses{Turing.Cfg, Turing.Cfg.mapState, Turing.CounterLoopState, Turing.MultiTapeTM, Turing.MultiTapeTM.runFrom, Turing.MultiTapeTM.runFrom_add, Turing.MultiTapeTM.runFrom_mapState_of_agreeOn, Turing.MultiTapeTM.step, Turing.counterLoopCfg, Turing.counterLoopRedirect, Turing.counterLoopTMVia, Turing.counterSpan, Turing.counter_layout, Turing.counter_redirect_injective, Turing.embedEmitCfg, Turing.embedEmitTM, Turing.embedEmitTM_runFrom}
\difficulty{6}
Let $k$ be a natural number, let $S$ and $\Gamma$ be types with decidable equality, and
let $e:\{\mathrm{false},\mathrm{true}\}\hookrightarrow\Gamma$ be an embedding. Let $M$
be a deterministic Turing machine over $\Gamma$ with $k$ work tapes and state set $S$,
and choose $a\in S$ and an optional exit state $\mathrm{exit}\in\mathrm{Option}(S)$. For
an input word $x$, let $c$ be a configuration of $M$ with state $a$, and let
$ct:\mathbb{Z}\to\mathrm{Option}(\Gamma)$ and $cp\in\mathbb{Z}$ specify the contents and
head position of an added counter tape; choose $t,\mathrm{width}\in\mathbb{N}$ with
$0<t$. Suppose that for every integer time $u$ with $0<u<t$, the run of $M$ from $c$ has
a state $s$ satisfying $s\ne a$ and $\mathrm{exit}\ne\mathrm{some}(s)$. Then after $t$
steps, the counter-loop machine formed from $M$, $e$, $a$, and $\mathrm{exit}$ has the
configuration obtained by embedding the configuration of $M$ after $t$ steps and
redirecting its optional state: states equal to $a$ enter the decrement phase, states
equal to $\mathrm{exit}$ enter the escape phase, and other states remain in the body
phase; the added counter tape and its head position are still $ct$ and $cp$. Moreover,
at every time $u<t$ the counter-loop machine is in neither its done nor escape state,
its added counter head lies in $[cp-1,cp+\mathrm{width}]$ at every time $u\le t$, and
there is some $u\le t$ at which all its original $k$ work-tape head positions agree with
those of $M$ after $u$ steps from $c$.
\end{lemma}

\begin{lemma}[Agreement of counter loops]
\lean{Turing.counter_refl_agree}
\label{Turing.counter_refl_agree}
\leanok
\uses{Turing.Action, Turing.Action.liftSymbols, Turing.Action.mapState, Turing.Action.mapWorkSymbols, Turing.FlagPhase, Turing.MultiTapeTM, Turing.MultiTapeTM.AgreeOn, Turing.MultiTapeTM.liftSymbols, Turing.MultiTapeTM.mapWorkSymbols, Turing.counterLoopDecExit, Turing.counterLoopTM, Turing.counterLoopTMVia, Turing.decodeBool, Turing.decrementTM, Turing.embedEmitTM, Turing.incrementTM, Turing.lastEmb}
\difficulty{6}
Let $k$ be a natural number, let $S$ be a state type with decidable equality, and let
$M$ be a deterministic $k$-work-tape machine over the Boolean alphabet with states in
$S$. For any $a\in S$ and optional exit state, the machine obtained by running the
counter loop over the Boolean alphabet and the machine obtained by running it through
the identity embedding of the Boolean alphabet into itself have identical transition
actions at every state, input symbol, and tuple of work-tape symbols.
\end{lemma}

\begin{lemma}[Counter loop rounds]
\lean{Turing.counter_rounds}
\label{Turing.counter_rounds}
\leanok
\uses{Turing.Cfg, Turing.FinTM.bufferTape, Turing.MultiTapeTM, Turing.MultiTapeTM.runFrom, Turing.MultiTapeTM.runFrom_add, Turing.MultiTapeTM.runFrom_zero, Turing.counterLoopCfg, Turing.counterLoopOrbit, Turing.counterLoopRedirect, Turing.counterLoopTMVia, Turing.counterOverhead, Turing.counterPatch, Turing.counterSpan, Turing.counterWord, Turing.counterWord_length, Turing.counter_body, Turing.counter_debit, Turing.counter_layout, Turing.counter_patch_frame, Turing.counter_patch_twice, Turing.counter_span_append, Turing.counter_word_decrement, Turing.decFixed, Turing.roundOrbit}
\difficulty{7}
Let $k$ be a natural number, let $\Gamma$ and $S$ be types with decidable equality, let
$x$ be a finite input word over $\Gamma$, let
$e:\{\mathrm{false},\mathrm{true}\}\hookrightarrow\Gamma$, and let $M$ be a
$k$-work-tape machine over $\Gamma$ with states in $S$. Fix a state $a\in S$, an
optional exit state, a time function $\tau$ on configurations, an initial configuration
$c$, a Boolean word $w$, a tape $ct:\mathbb Z\to\mathrm{Option}(\Gamma)$, a position
$cp\in\mathbb Z$, and a natural number $n$; suppose that for every integer $p$ with
$-1\le p\le |w|$, the tape cell $ct(cp+p)$ contains the symbol stored at offset $p$ in
the word $w$ with its bits mapped by $e$ (and is blank when that offset is outside the
word), and that $n$ is at most the natural number represented by the binary digits of
$w$. Suppose also that, for every $r<n$, the configuration after $r$ counter-loop rounds
has state $a$, its prescribed duration is positive, and at each strictly intermediate
positive time the machine is in a state different from $a$ and different from the
designated exit state. Then after the sum of those $n$ durations plus the counter
overhead, the counter-loop machine started with counter state in its initial decrement
phase is in the configuration obtained from the $n$-round orbit of $c$ by redirecting
its state (with the initial decrement phase used when $n=0$), and by replacing the
counter segment of $ct$ with the word obtained after $n$ fixed-width decrements, mapped
by $e$. Throughout this run, before the final time the machine has not reached either
exit anchor; at every time up to and including the final time, the counter head remains
between $cp-1$ and $cp+|w|$, inclusive, and the other work-head positions satisfy the
stated round-by-round condition: at each such time they equal those of the base machine
at some time within one of the first $n$ rounds, or at time zero of round $n$.
\end{lemma}

\begin{theorem}[Counter loop run to completion]
\lean{Turing.counterLoopTMVia_run_done}
\label{Turing.counterLoopTMVia_run_done}
\leanok
\uses{Turing.Cfg, Turing.FinTM.bufferTape, Turing.MultiTapeTM, Turing.MultiTapeTM.runFrom, Turing.MultiTapeTM.runFrom_add, Turing.counterLoopCfg, Turing.counterLoopOrbit, Turing.counterLoopRedirect, Turing.counterLoopTMVia, Turing.counterOverhead, Turing.counterPatch, Turing.counterWord, Turing.counterWord_length, Turing.counter_debit, Turing.counter_dec_zero, Turing.counter_patch_frame, Turing.counter_patch_twice, Turing.counter_rounds, Turing.counter_span_append, Turing.counter_word_exhausted, Turing.decFixed}
\difficulty{6}
Let $M$ be a deterministic machine with $k$ Boolean work tapes and state $a$, and let
$exit$ be an optional state. Let $w$ be a Boolean word with binary value $d$, and
suppose the extra tape initially agrees with the word buffer at cells $cp+p$ for all
integers $p$ with $-1\le p\le |w|$. Given an initial configuration $c$ and a
run-length function $\tau$, suppose that for every $r<d$, the $r$th round configuration
is in state $a$, its run length is positive, every intermediate configuration at a
positive time strictly before the round endpoint has a state different from $a$ and is
not the designated exit, and the next round configuration is again in state $a$. If $T$
equals the sum of these round lengths plus $\operatorname{overhead}(w,d+1)$, then the
counter-loop machine reaches the done configuration at time $T$, with the extra tape
filled with true symbols on cells $cp$ through $cp+|w|-1$. At all earlier times it is
neither done nor escape; through time $T$ the extra head lies in $[cp-1,cp+|w|]$, and
the first $k$ work-head positions at each time coincide with those in a run of $M$ from
some round configuration $c_r$, at a time $u$ satisfying either $r<d$ and
$u\le\tau(c_r)$ or $r=d$ and $u=0$.
\end{theorem}

\begin{theorem}[Counter loop run to escape]
\lean{Turing.counterLoopTMVia_run_escape}
\label{Turing.counterLoopTMVia_run_escape}
\leanok
\uses{Turing.Cfg, Turing.FinTM.bufferTape, Turing.MultiTapeTM, Turing.MultiTapeTM.runFrom, Turing.counterLoopCfg, Turing.counterLoopOrbit, Turing.counterLoopRedirect, Turing.counterLoopTMVia, Turing.counterOverhead, Turing.counterPatch, Turing.counterWord, Turing.counterWord_length, Turing.counter_rounds}
\difficulty{3}
Let $M$ be a deterministic machine with $k$ Boolean work tapes and distinct states $a$
and $e$. Let $w$ be a Boolean word, and suppose the extra tape initially agrees with the
word buffer at cells $cp+p$ for all integers $p$ with $-1\le p\le |w|$. Given an
initial configuration $c$, a run-length function $\tau$, and an index $r_0$ strictly
below the binary value of $w$, suppose each round $r\le r_0$ starts in state $a$, has
positive length, and at every positive intermediate time has a state different from both
$a$ and $e$. Suppose also that the configuration after round $r_0$ is in state $e$. If
$T$ is the sum of the round lengths for $r\le r_0$ plus
$\operatorname{overhead}(w,r_0+1)$, then the counter-loop machine reaches its escape
configuration at time $T$, with its extra tape on the word cells equal to the counter
word after $r_0+1$ decrements. Before $T$ the state is neither done nor escape; through
$T$ the extra head lies in $[cp-1,cp+|w|]$, and the first $k$ head positions at each
time agree with those of a run of $M$ from some round configuration, either with $r\le
r_0$ and $u\le\tau(c_r)$ or with $r=r_0+1$ and $u=0$.
\end{theorem}

\begin{theorem}[Boolean counter loop run to completion]
\lean{Turing.counterLoopTM_run_done}
\label{Turing.counterLoopTM_run_done}
\leanok
\uses{Turing.Cfg, Turing.CounterLoopState, Turing.FinTM.bufferTape, Turing.MultiTapeTM, Turing.MultiTapeTM.runFrom, Turing.MultiTapeTM.runFrom_eq_of_agreeOn, Turing.counterLoopCfg, Turing.counterLoopOrbit, Turing.counterLoopTM, Turing.counterLoopTMVia, Turing.counterLoopTMVia_run_done, Turing.counterOverhead, Turing.counter_refl_agree}
\difficulty{3}
Let $M$ be a deterministic machine with $k$ Boolean work tapes and state $a$, and let
$exit$ be an optional state. Let $w$ be a Boolean word with binary value $d$, and
suppose the extra tape initially agrees with the word buffer at cells $cp+p$ for every
integer $p$ with $-1\le p\le |w|$. Given an initial configuration $c$ and a run-length
function $\tau$, suppose each round $r<d$ starts in state $a$, has positive length, and
at every positive time strictly before its endpoint has a state different from $a$ and
not equal to the designated exit; also suppose each next round configuration is in state
$a$. If $T$ equals the sum of these round lengths plus
$\operatorname{overhead}(w,d+1)$, the Boolean counter-loop machine reaches its done
configuration at time $T$, with the extra tape set to true on cells $cp$ through
$cp+|w|-1$. Before time $T$ its state is neither done nor escape; through $T$ the extra
head remains in $[cp-1,cp+|w|]$, and the first $k$ work-head positions at each time
match a run of $M$ from some round configuration, with either $r<d$ and
$u\le\tau(c_r)$ or $r=d$ and $u=0$.
\end{theorem}

\begin{theorem}[Boolean counter loop run to escape]
\lean{Turing.counterLoopTM_run_escape}
\label{Turing.counterLoopTM_run_escape}
\leanok
\uses{Turing.Cfg, Turing.CounterLoopState, Turing.FinTM.bufferTape, Turing.MultiTapeTM, Turing.MultiTapeTM.runFrom, Turing.MultiTapeTM.runFrom_eq_of_agreeOn, Turing.counterLoopCfg, Turing.counterLoopOrbit, Turing.counterLoopTM, Turing.counterLoopTMVia, Turing.counterLoopTMVia_run_escape, Turing.counterOverhead, Turing.counterWord, Turing.counter_refl_agree}
\difficulty{3}
Let $M$ be a deterministic machine with $k$ Boolean work tapes and distinct states $a$
and $e$. Let $w$ be a Boolean word, and suppose the extra tape initially agrees with the
word buffer at cells $cp+p$ for every integer $p$ with $-1\le p\le |w|$. Given an
initial configuration $c$, a run-length function $\tau$, and an index $r_0$ strictly
below the binary value of $w$, suppose every round $r\le r_0$ starts in state $a$, has
positive length, and all positive-time intermediate states differ from both $a$ and $e$.
Suppose the configuration after round $r_0$ has state $e$. If $T$ is the sum of the
round lengths for $r\le r_0$ plus $\operatorname{overhead}(w,r_0+1)$, then the Boolean
counter-loop machine reaches its escape configuration at time $T$, with the extra tape
holding the counter word after $r_0+1$ decrements. Before $T$ its state is neither done
nor escape; through $T$ its extra head lies in $[cp-1,cp+|w|]$, and its first $k$ head
positions match those of some run of $M$ from a round configuration, with either $r\le
r_0$ and $u\le\tau(c_r)$ or $r=r_0+1$ and $u=0$.
\end{theorem}

\begin{theorem}[Time bound for the counter loop]
\lean{Turing.counterLoop_time_le}
\label{Turing.counterLoop_time_le}
\leanok
\uses{Turing.Cfg, Turing.MultiTapeTM, Turing.counterLoopOrbit, Turing.counterOverhead, Turing.counterOverhead_le}
\difficulty{3}
Let $M$ be a multi-tape Turing machine, let $x$ be an input word, let $c$ be a
configuration of $M$ on $x$, and let $\tau$ assign a natural-number time bound to each
such configuration. For a Boolean word $w$ whose binary value is $d$, and a natural
number $B$ such that $\tau$ is at most $B$ on each configuration reached after $r$
counter-loop rounds for every $r<d$, the sum of these $d$ time bounds plus the counter
overhead of $w$ through round $d$ is at most $dB+4d+2|w|+2$.
\end{theorem}

\begin{definition}[Counter loop expansion]
\lean{Turing.expandTM}
\label{Turing.expandTM}
\leanok
\uses{Turing.CounterLoopState, Turing.MultiTapeTM, Turing.counterLoopTM, Turing.writeAdvTM}
Given a number $k$ of work tapes, an index $i<k$, and a finite Boolean pattern,
construct a machine with $k+1$ work tapes and Boolean work-tape symbols. Its states are
the loop states carrying states of the machine that writes the pattern while advancing
on tape $i$; the constructed machine runs the counter loop around that writer, with the
counter initialized to $0$ and no designated exit state.
\end{definition}

\begin{theorem}[Expansion machine simulation]
\lean{Turing.expandTM_run_ofCfg}
\label{Turing.expandTM_run_ofCfg}
\leanok
\uses{Turing.Cfg, Turing.FinTM.bufferTape, Turing.MultiTapeTM.runFrom, Turing.MultiTapeTM.runFrom_state_ne_none_of_le, Turing.MultiTapeTM.runFrom_zero, Turing.counterLoopCfg, Turing.counterLoopOrbit, Turing.counterLoopTM_run_done, Turing.counterOverhead, Turing.expandTM, Turing.writeAdvTM, Turing.writeAdvTM_run_ofCfg}
\difficulty{6}
Let $k$ be a natural number, let $x$ and $w$ be Boolean words, let $i$ be one of $k$
work tapes, let $\mathrm{pat}$ be a Boolean word, and let $c$ be a configuration of a
$k$-tape Boolean machine whose state is $0$. Suppose a tape content $\mathrm{ct}$ agrees
with $w$ on the interval from $-1$ through $|w|$ when read relative to head position
$\mathrm{cp}$, and let $d=\operatorname{val}_2(w)$. Run the machine that repeatedly
writes the pattern on tape $i$, using the binary word as its counter, for
$d(|\mathrm{pat}|+1)+\mathrm{overhead}(w,d+1)$ steps from the counter-loop configuration
built from $c$, $\mathrm{ct}$, and $\mathrm{cp}$. The resulting configuration is the
counter-loop configuration in state done whose simulated configuration has tape $i$
patched with $d$ consecutive copies of $\mathrm{pat}$, whose head on tape $i$ has
advanced by $d|\mathrm{pat}|$, and whose counter tape is true exactly on the original
word interval; before that time it is in neither done nor escape. Throughout the run,
the counter head stays between $\mathrm{cp}-1$ and $\mathrm{cp}+|w|$, and each simulated
tape head stays between its starting position and its position advanced by
$d|\mathrm{pat}|$ on tape $i$, while every other simulated tape head remains fixed.
\end{theorem}

\begin{theorem}[Counter loop completion witness]
\lean{Turing.counterLoopTMVia_done_witness}
\label{Turing.counterLoopTMVia_done_witness}
\leanok
\uses{Turing.Cfg, Turing.FinTM.bufferTape, Turing.MultiTapeTM, Turing.MultiTapeTM.runFrom, Turing.Round, Turing.Round.cut_some, Turing.Segment, Turing.Segment.live, Turing.counterLoopCfg, Turing.counterLoopTMVia, Turing.counterLoopTMVia_run_done, Turing.counterOverhead, Turing.roundOrbit}
\difficulty{4}
Let $M$ be a deterministic machine with $k$ work tapes and state $a$, and let
$e:\{0,1\}\hookrightarrow\Gamma$ encode Boolean symbols. Let $w$ be a Boolean word
with binary value $d$, and suppose the extra tape initially agrees with the encoded word
buffer at every cell $cp+p$ for integers $p$ satisfying $-1\le p\le |w|$. For an
initial configuration $c$, a run-length function $\tau$, and a property $P$ of the
first $k$ work-head positions, suppose that every round $r<d$ satisfies the round
conditions for exit set $\{exit\}$ and property $P$, and that its next round
configuration is in state $a$. Then the counter-loop machine reaches its done
configuration after $\sum_{r<d}\tau(c_r)+\operatorname{overhead}(w,d+1)$ steps, where
$c_r$ is the $r$th round configuration and the extra tape contains encoded true symbols
on cells $cp$ through $cp+|w|-1$. Before that time the machine's state is never halting,
done, or escape. At every time through the endpoint, the extra head lies in
$[cp-1,cp+|w|]$, and there are round index $r$ and within-round time $u$ such that
either $r<d$ and $u\le\tau(c_r)$, or $r=d$ and $u=0$, with the first $k$ head
positions equal to those in the corresponding run of $M$.
\end{theorem}

\begin{theorem}[Counter loop completion segment]
\lean{Turing.counterLoopTMVia_done_segment}
\label{Turing.counterLoopTMVia_done_segment}
\leanok
\uses{Turing.Cfg, Turing.FinTM.bufferTape, Turing.MultiTapeTM, Turing.MultiTapeTM.runFrom, Turing.Round, Turing.Segment, Turing.Segment.mono, Turing.counterLoopCfg, Turing.counterLoopTMVia, Turing.counterLoopTMVia_done_witness, Turing.counterOverhead, Turing.roundOrbit}
\difficulty{5}
Let $M$ be a deterministic machine with $k$ work tapes, state $a$, designated exit state
or halting exit, and let $e:\{0,1\}\hookrightarrow\Gamma$ encode Boolean symbols in
its alphabet. Let $w$ be a Boolean word whose binary value is $d$, and suppose the extra
tape initially stores the encoded word at cells $cp$ through $cp+|w|-1$: for every
integer $p$ with $-1\le p\le |w|$, its content at $cp+p$ equals the word-buffer
content at $p$. Let $c$ be an initial configuration and let $P$ be a property of the $k$
work-head positions, with $P$ true at $c$. Suppose that for each $r<d$, the $r$th
iterate of the variable-length run of $M$ from $c$ to its next round has a round of
length $\tau$, satisfies the round conditions for exit set $\{exit\}$ and property
$P$, and ends in state $a$. Then the counter-loop machine that runs $M$ on the first $k$
tapes and uses the encoded extra tape for its counter reaches its done configuration
after $\sum_{r<d}\tau(c_r)+\operatorname{overhead}(w,d+1)$ steps, where $c_r$ is the
$r$th round configuration and the counter tape has been filled with encoded true symbols
on the cells $cp$ through $cp+|w|-1$. Before that time its state is neither done nor
escape; throughout the segment the extra head stays in $[cp-1,cp+|w|]$, and $P$ holds
for the first $k$ head positions. Here a segment means the stated endpoint is reached,
no state in $\{done,escape\}$ occurs earlier, and the position condition holds at
every time up to and including the endpoint.
\end{theorem}

\begin{theorem}[Counter loop escape witness]
\lean{Turing.counterLoopTMVia_escape_witness}
\label{Turing.counterLoopTMVia_escape_witness}
\leanok
\uses{Turing.Cfg, Turing.FinTM.bufferTape, Turing.MultiTapeTM, Turing.MultiTapeTM.runFrom, Turing.Round, Turing.Round.cut_some, Turing.Segment, Turing.Segment.live, Turing.counterLoopCfg, Turing.counterLoopTMVia, Turing.counterLoopTMVia_run_escape, Turing.counterOverhead, Turing.counterWord, Turing.roundOrbit}
\difficulty{5}
Let $M$ be a deterministic machine with $k$ work tapes and distinct states $a$ and $ex$,
and let $e:\{0,1\}\hookrightarrow\Gamma$ encode Boolean symbols. Let $w$ be a
Boolean word and suppose the extra tape initially agrees with its encoded word buffer at
cells $cp+p$ for every integer $p$ with $-1\le p\le |w|$. For an initial configuration
$c$, run-length function $\tau$, and property $P$ of the first $k$ work-head positions,
choose $r_0$ strictly below the binary value of $w$. Suppose every round $r\le r_0$
satisfies the round conditions for exit set $\{some\\ ex\}$ and property $P$, and the
configuration after round $r_0$ is in state $ex$. Then after $\sum_{r\le
r_0}\tau(c_r)+\operatorname{overhead}(w,r_0+1)$ steps, the counter-loop machine is in
its escape configuration, with the extra tape on the word cells holding the encoded
buffer of the counter word after $r_0+1$ decrements. Before the endpoint its state is
never halting, done, or escape; through the endpoint the extra head remains in
$[cp-1,cp+|w|]$, and at each time the first $k$ head positions match those of some run
of $M$ from a round configuration, with either $r\le r_0$ and $u\le\tau(c_r)$ or
$r=r_0+1$ and $u=0$.
\end{theorem}

\begin{theorem}[Counter loop escape segment]
\lean{Turing.counterLoopTMVia_escape_segment}
\label{Turing.counterLoopTMVia_escape_segment}
\leanok
\uses{Turing.Cfg, Turing.FinTM.bufferTape, Turing.MultiTapeTM, Turing.MultiTapeTM.runFrom, Turing.Round, Turing.Segment, Turing.Segment.mono, Turing.counterLoopCfg, Turing.counterLoopTMVia, Turing.counterLoopTMVia_escape_witness, Turing.counterOverhead, Turing.counterWord, Turing.roundOrbit}
\difficulty{4}
Let $M$ be a deterministic machine with $k$ work tapes, distinct states $a$ and $ex$,
and let $e:\{0,1\}\hookrightarrow\Gamma$ encode Boolean symbols. Let $w$ be a
Boolean word, and suppose the extra tape initially agrees with the encoded word buffer
at cells $cp+p$ for every integer $p$ with $-1\le p\le |w|$. For an initial
configuration $c$, a run-length function $\tau$, and a property $P$ of the first $k$
work-head positions, choose $r_0$ strictly below the binary value of $w$. Suppose each
round $r\le r_0$ satisfies the round conditions for exit set $\{some\\ ex\}$ and
property $P$, and the configuration after round $r_0$ is in state $ex$. Then the
counter-loop machine reaches its escape configuration after $\sum_{r\le
r_0}\tau(c_r)+\operatorname{overhead}(w,r_0+1)$ steps, where $c_r$ is the $r$th round
configuration; its extra tape on the word cells agrees with the encoded counter word
after $r_0+1$ decrements. Before that time its state is neither done nor escape.
Throughout the segment the extra head lies in $[cp-1,cp+|w|]$, and $P$ holds for the
first $k$ head positions at every time through the endpoint.
\end{theorem}

\begin{definition}[Extracting a counter loop body configuration]
\lean{Turing.counterLoopBody}
\label{Turing.counterLoopBody}
\leanok
\uses{Turing.Cfg, Turing.CounterLoopState}
Given an index $a$ into a state set and a configuration of a machine with $k+1$ work
tapes whose states are counter-loop states, form a configuration of the original
$k$-tape machine by setting its state to $a$ and retaining the input-head position,
output, and the contents and head positions of the first $k$ work tapes.
\end{definition}

\begin{theorem}[Recovering a counter-loop body configuration]
\lean{Turing.counterLoopCfg_body}
\label{Turing.counterLoopCfg_body}
\leanok
\uses{Turing.Cfg, Turing.CounterLoopState, Turing.counterLoopBody, Turing.counterLoopCfg, Turing.embedEmitCfg_eq}
\difficulty{4}
Let $k\in\mathbb{N}$, let $\Gamma$ be an alphabet, let $x$ be a word over $\Gamma$, let
$S$ be a state type, let $a\in S$, and let $d$ be a configuration with $k+1$ work tapes
and state type given by the counter-loop states over $S$. Form the body configuration
from $d$ using state $a$ and the first $k$ tapes, and form the counter-loop
configuration from that body configuration using $d$'s state and its last tape's
contents and head position. The resulting configuration equals $d$.
\end{theorem}

\begin{definition}[Binary counter value]
\lean{Turing.counterValue}
\label{Turing.counterValue}
\leanok
For a finite Boolean word $w$, define its counter value to be the natural number
represented in base two by the sequence obtained by replacing each Boolean digit by its
natural-number value.
\end{definition}

\begin{theorem}[Binary digits of a Boolean word]
\lean{Turing.foldr_bit_eq_ofDigits}
\label{Turing.foldr_bit_eq_ofDigits}
\leanok
\difficulty{3}
For every finite Boolean word $w$, folding its bits from the right using the binary
place-value operation, starting from zero, gives the natural number whose base-two
digits are the natural-number values of the bits of $w$, in their original order.
\end{theorem}

\begin{theorem}[Value of the binary digit word]
\lean{Turing.counterValue_bits}
\label{Turing.counterValue_bits}
\leanok
\uses{Turing.codeBitsNat_bits, Turing.counterValue, Turing.foldr_bit_eq_ofDigits}
\difficulty{3}
For every natural number $t$, the binary value of the binary digit word of $t$ is $t$.
\end{theorem}

\begin{definition}[Contract for counter completion]
\lean{Turing.RoundContract.counterLoopDone}
\label{Turing.RoundContract.counterLoopDone}
\leanok
\uses{Turing.Contract, Turing.Contract.ofSegment, Turing.FinTM.bufferTape, Turing.MultiTapeTM, Turing.RoundContract, Turing.RoundContract.hullHi, Turing.RoundContract.hullLo, Turing.RoundContract.orbit_anchor, Turing.RoundContract.orbit_round, Turing.RoundContract.orbit_start, Turing.Segment.mono, Turing.counterLoopBody, Turing.counterLoopCfg, Turing.counterLoopCfg_body, Turing.counterLoopTMVia, Turing.counterLoopTMVia_done_segment, Turing.counterOverhead, Turing.counterValue, Turing.headBox_lastCases, Turing.roundOrbit, Turing.roundOrbit_succ}
Given a round contract for a machine with anchor state $a$ and escape set $\{e\}$, and
an embedding of the Boolean symbols into the machine's alphabet, construct a contract
for the counter-loop machine starting in its decrement phase with the run flag. The data
consist of an original parameter and a Boolean counter word; the precondition requires
that the extra tape contain the encoded counter word on the specified interval and that
every round before the counter value satisfies the round precondition and ends at the
anchor. The exact running time is the sum of the round times for rounds indexed from $0$
to one less than the counter value, plus the counter overhead for one more than the
counter value; the endpoint is the completed loop configuration after that many rounds,
with the counter tape patched to all true symbols on the specified interval. The
head-position bounds are the round lower and upper bounds on the original tapes, and on
the counter tape they are respectively one cell before its starting position and the
starting position plus the counter-word length.
\end{definition}

\begin{definition}[Contract for counter-loop escape]
\lean{Turing.RoundContract.counterLoopEscape}
\label{Turing.RoundContract.counterLoopEscape}
\leanok
\uses{Turing.Contract, Turing.Contract.ofSegment, Turing.FinTM.bufferTape, Turing.MultiTapeTM, Turing.RoundContract, Turing.RoundContract.hullHi, Turing.RoundContract.hullLo, Turing.RoundContract.orbit_anchor, Turing.RoundContract.orbit_round, Turing.Segment.mono, Turing.counterLoopBody, Turing.counterLoopCfg, Turing.counterLoopCfg_body, Turing.counterLoopTMVia, Turing.counterLoopTMVia_escape_segment, Turing.counterOverhead, Turing.counterValue, Turing.counterWord, Turing.headBox_lastCases, Turing.roundOrbit, Turing.roundOrbit_succ}
Let $a$ and $e$ be distinct states, and suppose a round contract for a machine with
anchor $a$ has escape set $\{e\}$. Given an embedding of Boolean symbols into the
machine's alphabet, construct a contract for the counter-loop machine starting in its
decrement phase with the run flag. Its data consist of an original parameter, a Boolean
counter word, and a natural number $t$. The precondition requires the encoded counter
word on the extra tape, $t$ strictly less than its counter value, every round through
round $t$ to satisfy the round precondition, all rounds before $t$ to end at the anchor,
and round $t$ to end at $e$. The exact running time sums the round times for rounds $0$
through $t$, together with the counter overhead for $t+1$ rounds. The endpoint is the
escaped configuration after $t+1$ rounds with the extra tape set to the counter word
decremented $t+1$ times; the head-position bounds are the round bounds on the original
tapes, and on the extra tape they are one cell before its initial head position and that
position plus the counter-word length.
\end{definition}

\begin{theorem}[Head Trajectories of a Counter Loop]
\lean{Turing.RoundContract.counterLoopDone_traj}
\label{Turing.RoundContract.counterLoopDone_traj}
\leanok
\uses{Turing.Contract.Traj, Turing.MultiTapeTM, Turing.MultiTapeTM.runFrom, Turing.RoundContract, Turing.RoundContract.counterLoopDone, Turing.RoundContract.orbit_anchor, Turing.RoundContract.orbit_round, Turing.counterLoopBody, Turing.counterLoopCfg_body, Turing.counterLoopTMVia_done_witness, Turing.counterValue, Turing.roundOrbit, Turing.roundOrbit_succ}
\difficulty{5}
Let $k$ be a natural number, let $S$ and $\Gamma$ be types with decidable equality, let
$e:\mathrm{Bool}\hookrightarrow\Gamma$ be an injection, let $M$ be a
$k$-tape Turing machine, and let $a\in S$ and $\mathrm{exit}:\mathrm{Option}(S)$. Let
$R$ be a round contract for $M$ with anchor $a$ and escape set $\{\mathrm{exit}\}$, and
fix an input $x$. For $p=(i,w)$ and a configuration $d$ of the counter-loop machine,
suppose the last tape of $d$ stores the encoded word $w$ at offsets $q$ for $-1\le q\le
|w|$, as specified by the buffer convention, and that every round with index $r< v(w)$,
starting from the body configuration associated with $d$, satisfies $R$'s precondition
and has an endpoint in state $a$. Then for every time $q$ up to the counter-loop
contract's prescribed duration, there exist $r$ and $u$ such that either $r<v(w)$ and
$u$ is at most the duration of round $r$, or $r=v(w)$ and $u=0$, and the first $k$
work-tape head positions at time $q$ equal those of the run of $M$ from the start of
round $r$ at time $u$.
\end{theorem}

\begin{theorem}[Trajectory of a Counter-Loop Escape]
\lean{Turing.RoundContract.counterLoopEscape_traj}
\label{Turing.RoundContract.counterLoopEscape_traj}
\leanok
\uses{Turing.Contract.Traj, Turing.MultiTapeTM, Turing.MultiTapeTM.runFrom, Turing.RoundContract, Turing.RoundContract.counterLoopEscape, Turing.RoundContract.orbit_anchor, Turing.RoundContract.orbit_round, Turing.counterLoopBody, Turing.counterLoopCfg_body, Turing.counterLoopTMVia_escape_witness, Turing.roundOrbit, Turing.roundOrbit_succ}
\difficulty{5}
Let $k$ be a natural number, let $S$ and $\Gamma$ have decidable equality, let
$e:\{\mathrm{false},\mathrm{true}\}\hookrightarrow\Gamma$ be an injection, let $M$ be a
$k$-tape Turing machine, and let $a,ex\in S$ with $ex\ne a$. Let $R$ be a round contract
for $M$ with anchor $a$ and escape set $\{\mathrm{some}(ex)\}$, and fix an input $x$.
For an index $i$, a Boolean word $w$, a natural number $r_0<v(w)$, and a configuration
$d$ of the counter-loop machine, suppose the last tape of $d$ stores the encoded word
$w$ at offsets $q$ for $-1\le q\le |w|$, as specified by the buffer convention; every
round with index $r\le r_0$ satisfies $R$'s round condition; and the endpoint of round
$r_0$ has state $ex$. Then for every time $q$ up to the counter-loop contract's
prescribed duration, there exist $r$ and $u$ such that either $r\le r_0$ and $u$ is at
most the duration of round $r$, or $r=r_0+1$ and $u=0$, and the first $k$ work-tape head
positions at time $q$ equal those of the run of $M$ from the start of round $r$ at time
$u$.
\end{theorem}
